\documentclass[aspectratio=169]{beamer}
\usepackage[T2A]{fontenc} % Без этого с компилятором pdflatex не работает
\usepackage[utf8]{inputenc} % Без этого с компилятором pdflatex не работает
\usepackage[russian]{babel} % Без этого с компилятором pdflatex не работает
\usepackage{graphicx} % Для вставки изображений
\usepackage{listings}

\usetheme{Madrid}
\usecolortheme{seahorse}
\setbeamertemplate{navigation symbols}{}
\setbeamertemplate{footline}{%
	\leavevmode%
	\hbox{%
		\begin{beamercolorbox}[wd=.6\paperwidth,ht=2.5ex,dp=1ex,left]{palette primary}%
			\hspace{0.5em}Студенты АлтГТУ, группа 8 ПИ-51
		\end{beamercolorbox}%
		\begin{beamercolorbox}[wd=.4\paperwidth,ht=2.5ex,dp=1ex,right]{palette primary}%
			\insertshortdate\hspace{1em}\insertframenumber/\inserttotalframenumber\hspace{0.5em}
		\end{beamercolorbox}%
	}%
	\vskip0pt%
}

\lstset{
	basicstyle=\ttfamily\scriptsize,
	breaklines=true,
	frame=single,
	showstringspaces=false
}

\title{Оценка реализации устойчивой к сбоям real-time OS}
\subtitle{Применение принципов отказоустойчивости BEAM к существующей реализации POK на C}

\author{
	Студенты АлтГТУ, группа 8 ПИ-51\\
	Бельш А.Е.\\
	Воробьев К.В.\\
	Мэн Ч.\\
	Ошурков Н.А.
}

\date{2026 г.}

\begin{document}
	
	% -------------------------------------------------
	\begin{frame}{Постановка задачи}
		\begin{block}{Основная идея}
			В работе BEAM рассматривается не как замена RTOS и не как новый runtime, а как источник архитектурных принципов отказоустойчивости.
		\end{block}
		
		\vspace{0.3cm}
		
		\begin{itemize}
			\item Базовая реализация POK уже существует на C.
			\item Задача состоит в том, чтобы определить, какие принципы BEAM можно наложить на существующий код.
			\item Основной акцент делается на:
			\begin{itemize}
				\item супервизии;
				\item обработке fault-events;
				\item restart from known state;
				\item разграничении инкарнаций после восстановления.
			\end{itemize}
		\end{itemize}
	\end{frame}
	
	% -------------------------------------------------
	\begin{frame}{Основной научный источник}
		\begin{block}{Inria Research Report No. 9511, 2023}
			Giovanni Fabbretti, Ivan Lanese, Jean-Bernard Stefani\\
			\textit{A Behavioral Theory For Crash Failures and Erlang-style Recoveries In Distributed Systems}	
		\end{block}
		
		\begin{itemize}
			\item В статье рассматриваются crash failures и recovery в распределённых системах.
			\item Авторы вводят модель, близкую к Erlang-style recovery.
			\item Важные для нашей темы свойства:
			\begin{itemize}
				\item asynchronous communication;
				\item dynamic nodes and links;
				\item crash failures with recovery;
				\item imperfect knowledge;
				\item incarnation numbers.
			\end{itemize}
		\end{itemize}
	\end{frame}
	
	% -------------------------------------------------
	\begin{frame}{Что видно в кодовой базе POK}
		\begin{itemize}
			\item В проекте отдельно отмечены незавершённые части fault-management.
			\item В \texttt{error.c} обработка ошибки по умолчанию не образует полноценного механизма восстановления.
			\item В \texttt{partition.c} восстановление partition помечено как неполное.
			\item В \texttt{sched.c} модель планирования остаётся ограниченной.
		\end{itemize}
		
		\vspace{0.3cm}	
			Проблема состоит не в отсутствии базовой RTOS-функциональности, а в неполной архитектурной оформленности fault-management слоя.
	\end{frame}
	
	% -------------------------------------------------
	\begin{frame}[fragile]{Симптомы в репозитории POK}
		\begin{lstlisting}
			KERNEL
			- Implement the health monitor that is near of the ARINC653 definition
			- Implement FIFO scheduler
			- Implement queueing discipline
			
			LIBPOK
			- Implement fault callback function
		\end{lstlisting}
		
		\vspace{0.3cm}
		
		\begin{itemize}
			\item Сам репозиторий показывает необходимость доработки мониторинга и обработки отказов.
			\item Следовательно, развитие fault tolerance для POK обосновано и по состоянию кода.
		\end{itemize}
	\end{frame}
	
	% -------------------------------------------------
	\begin{frame}{Проблема 1: нет целостной модели health monitoring}

			Fault-management присутствует фрагментарно, но не как единый слой управления жизненным циклом компонентов.
		
		\begin{block}{Что даёт BEAM}
			В OTP supervision tree:
			\begin{itemize}
				\item workers выполняют работу;
				\item supervisors отслеживают workers;
				\item supervisor перезапускает дочерний компонент по заданной политике;
				\item стратегия рестарта задаётся явно.
			\end{itemize}
		\end{block}
	\end{frame}
	
	% -------------------------------------------------
	\begin{frame}[fragile]{Схематичный перенос идеи supervisor на C}
		\begin{lstlisting}[language=C]
			typedef struct {
				int id;
				int restart_count;
				int max_restarts;
				int state;
			} child_t;
			
			void supervisor_on_fault(child_t* c) {
				if (c->restart_count < c->max_restarts) {
					restart_child(c->id);
					c->restart_count++;
				} else {
					set_degraded_mode(c->id);
				}
			}
		\end{lstlisting}
		
		\vspace{0.2cm}
		
			Идея заключается не в переносе Erlang runtime, а в переносе политики:
			отказ наблюдается извне и обрабатывается по заранее заданным правилам.
	\end{frame}
	
	% -------------------------------------------------
	\begin{frame}{Проблема 2: ошибка фиксируется, но не становится архитектурным событием}
		\begin{itemize}
			\item В \texttt{error.c} ошибка может приводить к warning, fatal или возврату.
			\item Однако сама ошибка не всегда превращается в полноценный fault-event.
			\item Это затрудняет отделение обнаружения отказа от решения о восстановлении.
		\end{itemize}
		
		\vspace{0.3cm}
		
		\begin{block}{Что даёт BEAM}
			В Erlang отказ становится наблюдаемым через сигналы и сообщения, а логика восстановления переносится во внешний supervisory слой.
		\end{block}
	\end{frame}
	
	% -------------------------------------------------
	\begin{frame}[fragile]{Схематичный fault-event bus на C}
		\begin{lstlisting}[language=C]
			typedef struct {
				unsigned kind;
				unsigned partition;
				unsigned thread;
				unsigned addr;
			} fault_event_t;
			
			void fault_publish(fault_event_t ev) {
				queue_push(&fault_bus, &ev);
			}
		\end{lstlisting}
		
		\vspace{0.2cm}
		
		\begin{block}{Смысл}
			Ошибка становится не локальной процедурной веткой, а событием,
			которое может быть обработано несколькими подсистемами:
			supervisor, logger, safe-mode handler.
		\end{block}
	\end{frame}
	
	% -------------------------------------------------
	\begin{frame}{Проблема 3: restart без clean recovery}
		\begin{itemize}
			\item В \texttt{pok\_partition\_reinit()} прямо отмечено, что queueing/sampling ports нужно дополнительно сбрасывать.
			\item Следовательно, после восстановления возможно смешение старого и нового IPC-состояния.
			\item В RTOS это опасно, поскольку старые сообщения могут перейти в новую фазу работы.
		\end{itemize}
		
		\vspace{0.3cm}
		
		\begin{block}{Что даёт статья 2023 года}
			Использование incarnation numbers позволяет отличать старую инкарнацию узла от новой и безопасно отбрасывать устаревшие сообщения.
		\end{block}
	\end{frame}
	
	% -------------------------------------------------
	\begin{frame}[fragile]{Схематичный перенос incarnation numbers}
		\begin{lstlisting}[language=C]
			typedef struct {
				unsigned epoch;
				unsigned msg_type;
				unsigned char payload[32];
			} msg_t;
			
			int accept_msg(msg_t* m, unsigned current_epoch) {
				return m->epoch == current_epoch;
			}
		\end{lstlisting}
		
		\vspace{0.2cm}
		
		\begin{block}{Смысл}
			После restart увеличивается epoch.
			Старые сообщения не проходят проверку и не загрязняют новое состояние partition.
		\end{block}
	\end{frame}
	
	% -------------------------------------------------
	\begin{frame}{Проблема 4: слишком много общего состояния}
		\begin{itemize}
			\item В POK активно используются глобальные структуры partition- и scheduler-уровня.
			\item Для RTOS это естественно, но усложняет локализацию отказов и частичный restart.
			\item Сильная связанность приводит к тому, что сбой одного компонента затрагивает больше контекста, чем хотелось бы.
		\end{itemize}
		
		\vspace{0.3cm}
		
		\begin{block}{Что даёт BEAM}
			BEAM опирается на ownership состояния и взаимодействие через сообщения.
			Это уменьшает связанность между компонентами и упрощает восстановление.
		\end{block}
	\end{frame}
	
	% -------------------------------------------------
	\begin{frame}{Как адаптировать принцип ownership к POK}
		\begin{itemize}
			\item Не требуется убирать все глобальные структуры.
			\item Требуется переразделить ответственность между подсистемами:
			\begin{itemize}
				\item scheduler owns scheduling state;
				\item IPC manager owns IPC state;
				\item supervisor owns lifecycle state;
				\item error manager owns fault state.
			\end{itemize}
			\item Такой подход можно рассматривать как actor-like discipline поверх уже существующей C-реализации.
		\end{itemize}
	\end{frame}
	
	% -------------------------------------------------
	\begin{frame}{Проблема 5: частичные и неоднозначные сбои}
		\begin{itemize}
			\item Для fault-tolerant RTOS недостаточно учитывать только fail-stop.
			\item После сбоя узлы могут иметь неактуальное представление о состоянии друг друга.
			\item Для группы вычислителей нужны:
			\begin{itemize}
				\item heartbeat;
				\item local view;
				\item suspicion level;
				\item majority voting;
				\item controlled rejoin после recovery.
			\end{itemize}
		\end{itemize}
		
		\vspace{0.3cm}
		
		\begin{block}{Роль опорной статьи}
			В статье imperfect knowledge рассматривается как важная часть recovery-модели.
			Это делает её особенно полезной для fault-tolerant RTOS.
		\end{block}
	\end{frame}
	
	% -------------------------------------------------
	\begin{frame}{Границы применимости BEAM}
		\begin{block}{Что не предлагается}
			\begin{itemize}
				\item не предлагается заменить hard real-time scheduler на BEAM;
				\item не предлагается переносить весь runtime Erlang внутрь POK.
			\end{itemize}
		\end{block}
		
		\begin{block}{Что предлагается}
			Использовать принципы BEAM только в control plane отказоустойчивости:
			\begin{itemize}
				\item supervision;
				\item fault-events;
				\item bounded restart;
				\item degraded mode;
				\item controlled recovery;
				\item reincarnation after restart.
			\end{itemize}
		\end{block}
	\end{frame}
	
	% -------------------------------------------------
	\begin{frame}{Итоговое предложение для POK}
		\begin{enumerate}
			\item Ввести supervisor layer поверх thread/partition механизмов.
			\item Перейти от callbacks к fault-event bus.
			\item Использовать incarnation numbers для partition и IPC.
			\item Переразделить ownership состояния между подсистемами.
			\item Для группы узлов использовать heartbeat, local view и voting.
			\item Реализовать Erlang-style recovery layer поверх существующего C-кода.
		\end{enumerate}
	\end{frame}
	
	% -------------------------------------------------
	\begin{frame}{Вывод}
		\begin{itemize}
			\item Кодовая база POK уже показывает необходимость развития fault-management слоя.
			\item Опорная статья 2023 года даёт формальную модель Erlang-style recovery.
			\item Наиболее переносимы из BEAM:
			\begin{itemize}
				\item супервизия;
				\item событийная обработка отказов;
				\item restart from known state;
				\item incarnation numbers;
				\item ownership-ориентированная организация состояния.
			\end{itemize}
			\item Следовательно, наиболее реалистичный путь --- построение Erlang-style recovery layer поверх существующего кода POK на C.
		\end{itemize}
	\end{frame}

	
\end{document}